Let $V=T\mathbin{\dot\cup}R$ be the set of locations that can be visited,
where $T$ contains the mandatory locations and $R$ contains the optional ones.
Let $W$ be a disjoint set of demand locations. The routing graph is the
complete undirected graph $G=(V,E)$, with $E=\{\{i,j\}:i,j\in V,\ i<j\}$
and edge costs $c_e$. For every $w\in W$, let $N_w\subseteq V$ denote
its admissible covering locations. In the supplied instances, $N_w$ consists
of the visitable locations whose input distance to $w$ does not exceed the
coverage threshold. A nonnegative prize $p_i$ is collected when $i\in V$
is visited, and the required total prize is $P$. A demand location is
served when at least one location in $N_w$ is visited. Coverage itself
does not add a prize to the visited locations' total.

The prize-collecting covering tour problem seeks a minimum-cost simple cycle
that visits every location in $T$, covers every demand in $W$, and collects
at least $P$ units of prize. This is the problem addressed by the supplied
MDM-GRASP implementation and manuscript~\cite{climaco2026mdm}.
We assume $T\ne\varnothing$ and select a root $r\in T$, consistently with
the benchmark instances used here. No triangle inequality is needed for
the formulation below.

For $S\subseteq V$, let $\delta(S)$ be the set of edges with exactly one
endpoint in $S$, and write $x(F)=\sum_{e\in F}x_e$ for $F\subseteq E$.
Binary variables $y_i$ indicate visited locations, and binary variables
$x_e$ indicate used routing edges. The following degree and connectivity
formulation serves as a reference model:
\begin{align}
 \min\quad &\sum_{e\in E}c_e x_e \label{eq:pcctp-objective}\\
 \text{subject to}\quad
 &x(\delta(\{i\}))=2y_i && i\in V,\label{eq:pcctp-degree}\\
 &y_i=1 &&i\in T,\label{eq:pcctp-mandatory}\\
 &\sum_{i\in N_w}y_i\ge1 &&w\in W,\label{eq:pcctp-cover}\\
 &\sum_{i\in V}p_i y_i\ge P,\label{eq:pcctp-prize}\\
 &x(\delta(S))\ge2y_i
 &&\varnothing\ne S\subseteq V\setminus\{r\},\ i\in S,
 \label{eq:pcctp-connectivity}\\
 &\sum_{i\in V}y_i\ge3,\label{eq:pcctp-size}\\
 &x_e\in\{0,1\} &&e\in E,\\
 &y_i\in\{0,1\} &&i\in V.
\end{align}
Constraints~\eqref{eq:pcctp-degree} assign degree two to visited locations
and degree zero to unvisited ones. Constraints~\eqref{eq:pcctp-cover} and
\eqref{eq:pcctp-prize} enforce coverage and prize collection, respectively.
The rooted connectivity inequalities~\eqref{eq:pcctp-connectivity} exclude
selected components disconnected from $r$. Constraint~\eqref{eq:pcctp-size}
explicitly states the simple-cycle convention; for integer solutions it also
follows from the degree equations, binary edges, and the mandatory root.
The objective minimizes the total cost of the selected routing edges.
This reference formulation is not claimed as a new formulation.

To verify the integer interpretation, any feasible PCCTP cycle directly
satisfies the model. Conversely, the degree equations imply that the graph
induced by selected edges is a collection of vertex-disjoint simple cycles.
If a selected component does not contain $r$, choosing its vertex set as
$S$ in~\eqref{eq:pcctp-connectivity} gives a contradiction. Thus all selected
locations belong to one cycle containing the mandatory root. The remaining
constraints give the required visits, coverage, and prize total.

Fixing a subset of these routing or visit decisions defines a restricted
repair problem. An optimality certificate for that restricted problem is
local to its fixed decisions. By itself, it does not establish global optimality when the fixed decisions
exclude feasible solutions of the reference problem.
This distinction applies to the contracted repair experiments reported in
our research record.
